
\documentclass[12pt]{article}
\usepackage{it_hw}
\newcommand{\E}[1]{\mathify{\mathbb{E}\left[#1\right]}}
\input{macro}

\includeversion{problem}
\excludeversion{solution}

\begin{document}

\begin{center}
{\bf Information Theory and Its Application in Deep Learning: Midterm \begin{solution}Solutions\end{solution}} \\ \medskip
\begin{problem}Tuesday, April 14, 2026 \end{problem}
\end{center}

\begin{problem}
\noindent 
\begin{itemize}
\item You can use previous parts of a problem even if you did not solve them.
\item As throughout the course, entropy ($H$) and Mutual Information ($I$) are in bits.
\item $\log$ is taken in base 2. 
\item `prefix code' refers to a variable length code satisfying the prefix condition.  
\end{itemize}
\LeftFoot{Midterm}
\end{problem}

\begin{solution}
\LeftFoot{Midterm Solution}
\end{solution}

\begin{problem}

    \vskip .3in
    {\bf ID : \underline{\hspace{6cm}}}\\
    \vskip .3in
    {\bf Name : \underline{\hspace{5.5cm}}}\\
    ~
    \vskip 2in
    \begin{table}[h]
        \Large
        \begin{tabular}{c|c}
           1 &\textcolor{white}{aaaaaaaa}/~50\\
           2 &\textcolor{white}{aaaaaaaa}/~25 \\
           3 &\textcolor{white}{aaaaaaaa}/~25 \\
            \hline
           Total &  \textcolor{white}{aaaaaaaa}/~100 \\
        \end{tabular}
    \end{table}
    \newpage
    {\bf extra sheet}
    \newpage
\end{problem}

\begin{enumerate}
\item  {\bf Mix of Questions} {\it (50 points, 10 points each)}
\begin{enumerate}[1)]
% (a)
\item Here is a statement about pairwise independence and joint independence. Let $X, Y_1,$ and $Y_2$ be binary random variables. If $I(X;Y_1)=0$ and $I(X;Y_2)=0$, does it follow that $I(X;Y_1, Y_2)=0$? Prove or provide a counterexample.


% (b)
\item  Suppose a (non-stationary) Markov chain starts in one of $n$ states, necks down to $k < n$ states,
and then fans back to $m > k$ states.
Thus $X_1 - X_2 - X_3$ forms a Markov chain, where $X_1 \in \{1,2,\ldots, n\}$, $X_2 \in  \{1,2,\ldots, k\}$, $X_3 \in \{1,2,\ldots, m\}$,
and $p(x_1,x_2,x_3)  = p(x_1)p(x_2|x_1)p(x_3|x_2)$.
Show that the dependence of $X_1$ and $X_3$ is limited by the bottleneck by proving that $ I(X_1 ;X_3) \leq \log k.$

% (c)
\item The following claim is sometimes found in the literature: 
``It can be shown that the length of the binary Huffman codeword of a symbol $a_i$ with
probability $p_i$ is always less than or equal to $\lceil -\log_2 p_i\rceil$''
Is this claim true? Justify your answer 

% (d)
\item The \emph{World Series} is a seven-game series that terminates as soon as either team wins four games.  
Let \(X\) be the random variable representing the full outcome of the series between teams \(A\) and \(B\);  
for example, possible values of \(X\) include sequences such as \(\text{AAAA}\), \(\text{BABABAB}\), and \(\text{BBBAAAA}\).  
Assuming that: i) the two teams are equally matched (each game is won by \(A\) or \(B\) with probability \(1/2\)), and
    ii) all games are independent. Compute $H(X)$.

% (e)
\item Consider an i.i.d.\ information source $X^n$ over $\mathcal{X}$ with pmf $p(x)>0$ for all $x \in \mathcal{X}$. 
Let $x^n = (x_1, x_2, \dots, x_n)$ be a sequence of length $n$ drawn from this source, 
and
let $N(x|x^n)$ be the number of occurrences of symbol $x$ in $x^n$ (e.g., $N(x=a|x^n=abcaba)$ is 3).
Then, define a strong typical set:
\[
T_{\epsilon}^{(n)}(X) = \{x^n : \left| \frac{1}{n} N(x|x^n) - p(x) \right| < \epsilon \mbox{ for all $x\in\mathcal{X}$}\}
\]
If a sequence $x^n$ is strongly $\epsilon$-typical ($x^n \in T_\epsilon^{(n)}(X)$), show that
\[
\left| -\frac{1}{n} \log p(x^n) - H(X) \right| < \frac{\epsilon}{p_{\min}} H(X),
\]
where $p_{\min} = \min_{x\in\mathcal{X}} p_X(x)$.

\end{enumerate}

%%% Solution of  Problem 1
%%%%%%%%%%%%%%%%%%%%%
\vskip 0.2in
\begin{solution}
{\bf Solution:}\\

\begin{enumerate}[1)]
%%(a) sol
\item The answer is ``no''.  Although at first the conjecture seems reasonable enough--after all, if $Y_1$ gives you no information about $X$, and if $Y_2$ gives you no information about $X$, then why should
the two of them together give any information?  But remember, it is NOT the case that $I(X;Y_1,Y_2) = I(X;Y_1) + I(X;Y_2)$.  The chain rule for information says instead that $I(X;Y_1, Y_2) = I(X;Y_1) + I(X;Y_2|Y_1)$.  The chain rule gives us reason to be skeptical about the conjecture.

This problem is reminiscent of the well-known fact in probability that pair-wise independence of three random variables is not sufficient to guarantee that all three are mutually independent.  $I(X;Y_1) = 0$ is equivalent to saying that $X$ and $Y_1$ are independent.  Similarly for $X$ and $Y_2$.  But just because $X$ is pairwise independent with each of $Y_1$ and $Y_2$, it does not follow that $X$ is independent of the vector $(Y_1,Y_2)$.

Here is a simple counterexample.  Let $Y_1$ and $Y_2$ be independent fair coin flips.  And let $X = Y_1$ XOR $Y_2$. $X$ is pairwise independent of both $Y_1$ and $Y_2$, but obviously not independent of the vector $(Y_1, Y_2)$, since $X$ is uniquely determined once you know $(Y_1, Y_2)$.

%% (b) sol
\item 
From the data processing inequality, and the fact that entropy is maximum for a uniform distribution, we get 
\begin{align*}
I(X_1;X_3)\leq& I(X_1;X_2) \\
=&  H(X_2)-H(X_2\bigm | X_1) \\
\leq&  H(X_2)           \\
\leq&  \log k. 
\end{align*}

%% (c) sol
\item This claim is not true in general! To prove it, we shall construct a counterexample.
Choose a small number $\epsilon$. Consider the following source $X$ with 4 symbols:
$$X =
\begin{cases} 
a_1, & \text{w.p. } p_1 = \epsilon \\ 
a_2, & \text{w.p. } p_2 = \frac{1}{3} - \epsilon \\ 
a_3, & \text{w.p. } p_3 = \frac{1}{3} - \epsilon \\ 
a_4, & \text{w.p. } p_4 = \frac{1}{3} + \epsilon 
\end{cases}$$

Now consider the Huffman tree corresponding to the above source. If we combine $a_1$ and $a_2$ in the first step, the code will have length functions: $l(a_1) = l(a_2) = 3$, $l(a_3) = 2$ and $l(a_4) = 1$. Now for $a_2$, the Shannon code length would be $\left\lceil \log_2 \frac{1}{\frac{1}{3} - \epsilon} \right\rceil = 2$, for $\epsilon < \frac{1}{12}$.

Thus, the binary Huffman codeword length can exceed the Shannon codeword length in general for a symbol.

%% (d) sol
\item 
Let \(Y\) denote the number of games played, so \(Y \in \{4,5,6,7\}\).  
The series lasts \(y\) games if the winning team wins the last game and has exactly 3 wins in the first \(y-1\) games.  
Thus
\[
P(Y=y) = 2 \times \frac{\binom{y-1}{3}}{2^y}, \qquad y=4,5,6,7,
\]
where the factor 2 accounts for symmetry between \(A\) and \(B\).

\[
\begin{array}{c|cccc}
y & 4 & 5 & 6 & 7 \\ \hline
P(Y=y) & 1/8 & 1/4 & 5/16 & 5/16
\end{array}
\]


Each possible outcome sequence \(x\) has probability \(P(x) = (1/2)^{|x|}\),  
and the number of valid sequences of length \(y\) is \(N_y = 2\binom{y-1}{3}\).  
Hence
\[
H(X) = -\sum_x P(x)\log_2 P(x)
     = \sum_{y=4}^{7} N_y (1/2)^y \, y
     = \sum_{y=4}^{7} y\,P(Y=y)
     = \frac{93}{16}~\text{bits.}
\]

%% (e) sol
\item 
Substitute the given formulas for the log-probability and the true entropy $H(X)$:
\begin{align*}
\left| -\frac{1}{n} \log p(x^n) - H(X) \right| 
&= \left| \sum_{x \in \mathcal{X}} \frac{1}{n} N(x|x^n)\log \frac{1}{p(x)} - \sum_{x \in \mathcal{X}} p(x) \log \frac{1}{p(x)} \right|\\
&= \left| \sum_{x \in \mathcal{X}} \left(\frac{1}{n} N(x|x^n) - p(x)\right)\log \frac{1}{p(x)}\right|\\
&\leq  \sum_{x \in \mathcal{X}} \left|\frac{1}{n} N(x|x^n) - p(x)\right|\log \frac{1}{p(x)}\\
&\leq  \sum_{x \in \mathcal{X}} \epsilon \log \frac{1}{p(x)}\\
&\leq  \sum_{x \in \mathcal{X}} \frac{\epsilon}{p_{\min}} p(x)\log \frac{1}{p(x)}\\
&= \frac{\epsilon}{p_{\min}} H(X).
\end{align*}
\end{enumerate}
\end{solution}



\begin{problem}
\newpage {\bf problem 1 cont.}
\end{problem}
\begin{problem}
\newpage
\end{problem}

%%% Problem 2: Bad Wine
%%%%%%%%%%%%%%%%%%%%%
\vskip 0.3in
\item {\bf Bad wine} {\it (25 points)}\\
There are 6 bottles of wine.
It is known that precisely one bottle has gone bad (tastes very bad).
From inspection of the bottles, it is determined that the probability $p_i$ that the $i^{\mbox{\footnotesize{th}}}$ bottle is bad is given by $(p_1, p_2, \ldots , p_6)=(\frac{8}{23}, \frac{6}{23}, \frac{4}{23}, \frac{2}{23}, \frac{2}{23}, \frac{1}{23})$.
Tasting will determine the bad wine.

Suppose you taste the wines one at a time.
Choose the order of tasting to minimize the expected number of tastings required to determine the bad bottle.
Remember, if the first 5 wines pass the test, you don't have to taste the last.

\begin{enumerate}
\item What is the expected number of tastings required? {\it (5 points)}
\item Which bottle should be tasted first? {\it (5 points)}
\end{enumerate}

Now, suppose you adopt a more sophisticated, \textbf{adaptive mixing strategy}. For any given tasting, you are allowed to pour a small sample from any chosen subset of the bottles into a fresh glass and taste the mixture. If the mixture tastes bad, the bad bottle is guaranteed to be within that subset. If it tastes fine, the bad bottle must be in the complementary subset. 

Furthermore, this strategy can be fully adaptive: your choice of which bottles to mix for the next tasting can depend entirely on the results of all previous tastings. You proceed in this manner, dynamically adjusting your mixtures based on prior results, and stop once the bad bottle is uniquely determined.

\begin{itemize}
\item[(c)] Show that designing this optimal adaptive mixing strategy is equivalent to designing an optimal prefix code for lossless data compression. {\it (10 points)}
\item[(d)] What is the minimum expected number of tastings required to determine the bad wine? What mixture should be tasted first? {\it (5 points)}
\end{itemize}




%%% Solution of  Problem 2: Bad Wine
%%%%%%%%%%%%%%%%%%%%%
\vskip 0.2 in
\begin{solution}
{\bf Solution: }\\
\begin{enumerate}
\item
If we taste one bottle at a time, the corresponding number of tastings are $\{1,2,3,4,5,5\}$ with some order. By the same argument as in Lemma 5.8.1, to minimize the expected length $\sum p_i l_k$ we should have $l_j \le l_k$ if $p_j > p_k$.  Hence, the best order of tasting should be from the most likely wine to be bad to the least. 

The expected number of tastings required is 
\begin{eqnarray*}
  \sum_{i=1}^{6}p_{i}l_{i} & = & 1 \times \frac{8}{23} + 
           2 \times \frac{6}{23} + 3 \times \frac{4}{23} +
           4 \times \frac{2}{23} + 5 \times \frac{2}{23} +
           5 \times \frac{1}{23} \\
                           & = & \frac{55}{23} \\
                           & \approx & 2.39
\end{eqnarray*}

\item   The first bottle to be tasted should be the one with probability $\frac{8}{23}$.

\item  In this adaptive mixing strategy, every tasting partitions the remaining possible bad bottles into two subsets (the bottles in the mixture vs. the bottles left out). The outcome of the tasting provides a binary answer (Yes/No: is the bad wine in the mixture?). 

This process directly corresponds to traversing a binary tree from the root to a leaf, where each leaf represents one of the 6 bottles. Assigning a bit $0$ or $1$ to each branch of the tree creates a prefix-free binary code for the 6 bottles. The number of tastings required to identify bottle $i$ is exactly the length $l_i$ of its corresponding codeword. Therefore, minimizing the expected number of tastings, $\sum p_i l_i$, is mathematically identical to minimizing the expected codeword length in a lossless compression scheme, which is achieved via Huffman coding.

\item   The Huffman coding is optimal.
\begin{center}
\[
\begin{array}{lllllllllllll}
(00)   && 8 && 8 && 8 && 9 && 14 && 23\\
(01)   && 6 && 6 && 6 && 8 && 9 \\
(11)   && 4 && 4 && 5 && 6\\
(101)  && 2 && 3 && 4 \\
(1000) && 2 && 2 \\
(1001) && 1 && 
\end{array}
\]
\end{center}

The expected number of tastings required is
\begin{eqnarray*}
   \sum_{i=1}^{6}p_{i}l_{i} & = & 2 \times \frac{8}{23} + 2 \times \frac{6}{23} + 2 \times \frac{4}{23} + 3 \times \frac{2}{23} + 4 \times \frac{2}{23} + 4 \times \frac{1}{23} \\
                            & = & \frac{54}{23} \\
                            & \approx & 2.35
\end{eqnarray*}
The mixture of the first and second bottles should be tasted first (or equivalently, the mixture of the third, fourth, fifth, and sixth). This splits the probability mass as evenly as possible based on the first branching of the Huffman tree.

\end{enumerate}

\end{solution}

\begin{problem}
\clearpage{\bf problem 2 cont.}
    ~
    \clearpage
\end{problem}


\item  {\bf JS-Divergence} {\it (25 points)}\\
Let $P$ and $Q$ be discrete probabilities on the alphabet $\{1,2,\cdots,K\}$.
Relative entropy $D(P\|Q)$, can be viewed as a distance between probability distributions.
In general it is not symmetric, i.e., there exist $P$ and $Q$ for which  $D(P\|Q) \neq D(Q\|P)$.
We can construct symmetric distance measures between two distributions based on relative entropy;
one such measure is called the Jensen-Shannon Divergence. 
For $R = (P+Q)/2$, i.e., $R$ is a new probability distribution on the alphabet $\{1, 2, \cdots, K\}$ such that
\[
R(i) = \frac{P(i)+Q(i)}{2} \quad\mbox{for all $i\in\{1, \dots, K\}$}
\]
the JS-Divergence is defined as
\[ 
JS(P\|Q) = \frac{1}{2}D \left( P \| R  \right)+ \frac{1}{2}D\left( Q\|R \right).
\] 
It is easy to check that $JS(P\|Q) = JS(Q\|P)$.

\begin{enumerate}
\item Let $Z$ be a Bernoulli($\frac{1}{2}$) random variable. Let $X$ be a random variable such that $X|\{Z=1\}$ has a distribution $P$
while $X|\{Z=0\}$ has a distribution $Q$.
Find the pmf $p_X$ of random variable $X$ as a function of $P, Q, R$. {\it (5 points)}

\item For the same $X$ and $Z$ in the previous problem, express $JS(P\|Q)$ in terms of $I(X;Z)$. {\it (10 points)}

\item What is the maximum value of $JS(P\|Q)$ among all possible distributions $P$ and $Q$. [Hint: Use the previous results.] {\it (10 points)}
\end{enumerate}


%%% Solution of  Problem
%%%%%%%%%%%%%%%%%%%%%
\vskip 0.2in
\begin{solution}
{\bf Solution:}\\

\begin{enumerate}
\item We have
\begin{align*}
    p_X(i) =& Pr(X=i)\\ 
    =& Pr(X=i|Z=1)Pr(Z=1) + Pr(X=i|Z=0)Pr(Z=0)\\
    =& \frac{1}{2}P(i) + \frac{1}{2}Q(i)\\
    =& R(i).
\end{align*}
Thus, $p_X = R$.
\item Correct Answer is $I(X;Z) = JS(P\|Q)$.\\Let $M = \frac{P+Q}{2}$. Note that X is distributed according to $M$ by definition. 
\begin{align}
I(X;Z) &= H(X)-H(X|Z)\\
&=\sum -M_i \log M_i - \frac{1}{2}\left(\sum -P_i\log P_i + \sum -Q_i\log Q_i\right)\\
&=\sum -\frac{P_i}{2} \log M_i + \sum -\frac{Q_i}{2} \log M_i - \frac{1}{2}\left(\sum -P_i\log P_i + \sum -Q_i\log Q_i\right)\\
&=\sum \frac{P_i}{2} (\log P_i - \log M_i) + \sum -\frac{Q_i}{2} (\log Q_i - \log M_i)\\
&=\frac{1}{2}\sum P_i\log \frac{P_i}{M_i} + \frac{1}{2}\sum Q_i\log\frac{Q_i}{M_i}\\
&=\frac{1}{2}D(P\|M)+\frac{1}{2}D(Q\|M)\\
&=JS(P\|Q)
\end{align}
An alternate proof, which exploits the representation of mutual information in terms of relative entropy is as follows:
\begin{align*}
I(X;Z) \stackrel{(a)}{=}& D(P_{X,Z} \| P_{X}P_{Z}) \\
=& D(P_{X|Z} \| P_{X}) \\
=& E \big[\log \frac{P_{X|Z}}{P_X}\big] \\
\stackrel{(b)}{=}& E \bigg[ E \big[\log \frac{P_{X|Z=z}}{P_X} | Z=z\big] \bigg]  \\
\stackrel{(c)}{=}& \frac{1}{2} D(P \| M) + \frac{1}{2} D(Q \| M) \\
=& JS(P\|Q). 
\end{align*}
In the above, (a) follows by definition of mutual information. (b) follows by tower property of expectation. (c) follows from the fact that $P_{X|Z=1} = P$ and $P_{X|Z=0} = Q$, and $P_{Z}(0) = P_{Z}(1) = \frac{1}{2}$. 
\item Correct answer is 1.\\Recall from part 1, that 
\[
JS(P\|Q) = I(X;Z) = H(Z)-H(Z|X) \leq H(Z) = 1.
\]
\end{enumerate}

\end{solution}

\begin{problem}
\newpage {\bf problem 3 cont.}
\end{problem}
\begin{problem}
\newpage{\bf extrasheet.}
\end{problem}
\begin{problem}
\newpage{\bf extrasheet.}
\end{problem}
\end{enumerate}


\end{document}